\documentclass{situationist-preprint}

\title{Davidson's Thesis on Metaphor Defended}
\msid{LNG-09-a}
\msfamily{LNG-09}
\msdiscipline{Philosophy of Language}
\mstype{Principal manuscript}
\mskeywords{metaphor; Davidson; literal meaning; speaker meaning; implicature; seeing-as; paraphrase}

\begin{document}
\maketitle

\begin{abstract}
When Romeo says that Juliet is the sun, does his sentence have, besides its literal meaning, a metaphorical meaning? Davidson argued that it does not: metaphors mean what the words, in their most literal interpretation, mean, and nothing more; what is distinctive of metaphor belongs to use. Most subsequent theories have rejected the thesis, locating a metaphorical content either in the sentence relative to context or in what the speaker means. We defend Davidson's thesis. We argue that a metaphorical meaning would have to be fixed either by convention or by a rule of interpretation applying to novel metaphors, and that neither is available: conventions do not cover novel metaphors, and the effects of a metaphor are open-ended in a way no rule fixes. What speakers sometimes convey by a metaphor is a conversational implicature, which belongs to use; and what metaphor distinctively does, getting a hearer to see one thing as another, is not a content at all. The thesis survives the objection from metaphorical truth and the objection from disagreement about interpretations.
\end{abstract}

\section{The problem}

\citet{Black1954} held that a metaphor creates a new meaning through the interaction of two systems of associated commonplaces, and that its cognitive content cannot be captured by any literal paraphrase. \citet{Searle1979} located metaphorical meaning in speaker meaning: the speaker says that $S$ is $P$ and means that $S$ is $R$, and principles of inference take the hearer from one to the other. \citet{Stern2000} treats metaphorical interpretation as the evaluation of a context-sensitive operator, so that metaphors have a semantic content relative to context.

\citet{Davidson1978} rejected all such views. Metaphors mean what the words, in their most literal interpretation, mean, and nothing more. What a metaphor makes us notice is neither finite in scope nor propositional in character; to ascribe it a content is to confuse what it means with what it makes us see. We defend the thesis, and add an account of how metaphors can be used to convey determinate propositions.

\section{Framework}

\begin{definition}[Metaphorical meaning]\label{d:mm}
A sentence used metaphorically has a \emph{metaphorical meaning} if there is a proposition, other than its literal content, that the sentence expresses on that use in virtue of semantic conventions or rules of the language, together with the context.
\end{definition}

\begin{principle}[Meaning requires rules]\label{p:rules}
Whatever a sentence means, in a language, on a use, it means in virtue of conventions or rules that a competent speaker knows, and that determine its content from the meanings of its parts and the context.
\end{principle}

\begin{principle}[Use]\label{p:use}
A speaker may, by uttering a sentence with a given meaning, bring about effects in a hearer, including beliefs, images and changes of attention, that are not part of what the sentence means. Among them are conversational implicatures, which are calculable from the maxims and cancellable (LNG-11).
\end{principle}

\section{The argument}

\begin{lemma}\label{l:novel}
Novel metaphors have no metaphorical meaning in virtue of convention.
\end{lemma}

\begin{proof}
A metaphor is novel if its use is not the repetition of an established use. No convention covers what has not been established, so no convention assigns a novel metaphor a content beyond its literal one. When a metaphor becomes conventional, as \enquote{mouth of the river} has, the expression acquires a new literal meaning, and the metaphor is dead; this is Davidson's point that dead metaphors are not metaphors.
\end{proof}

\begin{lemma}\label{l:rules}
There is no rule that determines, for a novel metaphor in context, a proposition it expresses.
\end{lemma}

\begin{proof}
Suppose there were such a rule. Then a competent speaker who knows the rule and the context would know, for any novel metaphor, which proposition it expresses, as she knows what \enquote{I am hungry} expresses when she knows who speaks. But competent and attentive hearers diverge in what they take a novel metaphor to convey, and, when they agree on some proposition it may be used to convey, they continue to find more in it, without any point at which the content is fixed. \enquote{Juliet is the sun} may convey that she is radiant, that she is the centre of Romeo's world, that his day begins with her, that she sustains him, and the list does not close. A rule determining content in context would yield a determinate proposition, or at least a determinate range; the open-endedness of interpretation shows that none does.
\end{proof}

\begin{theorem}\label{t:main}
Metaphors mean exactly what they literally mean. Their distinctive effect lies in use.
\end{theorem}

\begin{proof}
By Principle~\ref{p:rules}, a metaphorical meaning (Definition~\ref{d:mm}) would have to be fixed by conventions or by rules applying in context. By Lemma~\ref{l:novel}, conventions do not fix one for novel metaphors, and conventional metaphors have only a literal meaning, the new one they have acquired. By Lemma~\ref{l:rules}, no rule fixes one either. So metaphorical sentences have no meaning beyond their literal meaning. Their effects in a hearer, which are what makes them metaphors, are effects of uttering a sentence with that meaning, and by Principle~\ref{p:use} they belong to use.
\end{proof}

The theorem does not deny that speakers can communicate determinate propositions by metaphor. A speaker who says, of a colleague who never changes his mind, \enquote{He is a rock} may intend to convey that he is unyielding, and the hearer may recover this by reasoning from the patent falsity of the literal content and the presumption of cooperation, as \citet{Grice1975} proposed. That proposition is a conversational implicature and, by Principle~\ref{p:use}, belongs to use. It is not the metaphor's meaning, and it does not exhaust what the metaphor does: having grasped that he is unyielding, the hearer has not yet seen him as a rock. What metaphor distinctively does is to get the hearer to see one thing as another, and seeing-as is not a content \citep{Davidson1978,Moran1989}.

\section{Objections and replies}

\begin{objection}[Metaphorical truth]\label{o:truth}
We assess metaphors as true or false. \enquote{Juliet is the sun} is not true, but \enquote{Man is a wolf to man} may be. If metaphors only meant what they literally mean, almost all would be false, and the assessment would be inexplicable.
\end{objection}

\begin{reply}
Most metaphors are literally false, and we know it; that is often how we recognize them. When we call a metaphor \enquote{true}, we assess what it is used to convey or what it makes us notice, as we call a joke apt. The relation between metaphor and literal falsity is examined in the companion manuscript LNG-09-b.
\end{reply}

\begin{objection}[Paraphrase]\label{o:para}
Metaphors can be paraphrased, at least partially, and teachers, critics and readers agree on many paraphrases. If there were nothing to paraphrase, this agreement would be inexplicable.
\end{objection}

\begin{reply}
Agreement on what a metaphor is used to convey is agreement about use, explained by shared background and shared reasoning about the speaker's purposes. The partial character of paraphrase, which \citet{Black1954} stressed, is what our account predicts: a paraphrase states some of the things the metaphor makes us notice, and the rest is not a remainder of content but a further effect \citep{Camp2006}.
\end{reply}

\begin{objection}[Embedding]\label{o:embed}
Metaphors embed under negation, conditionals and attitude verbs: \enquote{If Juliet is the sun, Rosaline is a candle}. Implicatures do not embed in this way, so metaphorical content must be part of what is said \citep{Stern2000}.
\end{objection}

\begin{reply}
The conditional is itself used metaphorically, as a whole; it means what it literally means, and what it is used to convey is recovered by reasoning about the whole utterance. \citet{LeporeStone2010} argue on independent grounds that the evidence from embedding does not establish metaphorical meaning.
\end{reply}

\section{Conclusion}

A metaphorical sentence means what it literally means. A metaphorical meaning would have to be fixed by convention or by rule, and novel metaphors have neither; dead metaphors, which have conventions, have a new literal meaning instead. What speakers convey by metaphor, when it is determinate, is a conversational implicature, and what metaphor distinctively does, making us see one thing as another, is an effect of use, not a meaning.

\begin{thebibliography}{99}\small
\bibitem[Black(1954)]{Black1954} Black, M. (1954). Metaphor. \emph{Proceedings of the Aristotelian Society}, 55, 273--294.
\bibitem[Camp(2006)]{Camp2006} Camp, E. (2006). Metaphor and that certain \enquote{je ne sais quoi}. \emph{Philosophical Studies}, 129(1), 1--25.
\bibitem[Davidson(1978)]{Davidson1978} Davidson, D. (1978). What metaphors mean. \emph{Critical Inquiry}, 5(1), 31--47.
\bibitem[Grice(1975)]{Grice1975} Grice, H.~P. (1975). Logic and conversation. In P. Cole and J.~L. Morgan (eds), \emph{Syntax and Semantics 3: Speech Acts}, 41--58. Academic Press.
\bibitem[Lepore and Stone(2010)]{LeporeStone2010} Lepore, E. and Stone, M. (2010). Against metaphorical meaning. \emph{Topoi}, 29(2), 165--180.
\bibitem[Moran(1989)]{Moran1989} Moran, R. (1989). Seeing and believing: Metaphor, image, and force. \emph{Critical Inquiry}, 16(1).
\bibitem[Searle(1979)]{Searle1979} Searle, J.~R. (1979). Metaphor. In A. Ortony (ed.), \emph{Metaphor and Thought}. Cambridge University Press.
\bibitem[Stern(2000)]{Stern2000} Stern, J. (2000). \emph{Metaphor in Context}. MIT Press.
\end{thebibliography}

\end{document}
